\documentclass[11pt]{article}
\usepackage[margin=1in]{geometry}
\usepackage{enumitem}
% =============================================================================
% Preambles/header.tex — shared LaTeX/Beamer preamble for this project
%
% Use from a Beamer lecture by prepending:
%     \input{header}
% and compiling with TEXINPUTS=../Preambles:$TEXINPUTS (the /compile-latex
% skill does this automatically).
%
% PALETTE CONTRACT. Color names defined here MUST match the names in
% Quarto/theme-template.scss so Beamer and Quarto renderings use the same
% palette. The scripts/check-palette-sync.sh script enforces this.
%
% To customize: edit the HEX values below AND the matching $variable in
% Quarto/theme-template.scss, then run scripts/check-palette-sync.sh.
% =============================================================================

% -----------------------------------------------------------------------------
% Engine detection + encoding
%
% Our /compile-latex skill uses XeLaTeX, where inputenc is unnecessary and
% can produce encoding warnings. Only load inputenc under pdfTeX.
% -----------------------------------------------------------------------------
\usepackage{iftex}
\ifPDFTeX
  \usepackage[utf8]{inputenc}
\fi

\usepackage{amsmath, amssymb, amsthm}
\usepackage{booktabs}
\usepackage{graphicx}

% -----------------------------------------------------------------------------
% PALETTE — mirrors Quarto/theme-template.scss variable names.
% Load xcolor and define colors BEFORE anything that references them
% (notably hyperref, hypersetup, and the Beamer color assignments).
% -----------------------------------------------------------------------------
\usepackage{xcolor}

\definecolor{primary-blue}{HTML}{012169}      % headings, accents
\definecolor{primary-gold}{HTML}{B9975B}      % emphasis, borders
\definecolor{highlight-yellow}{HTML}{F2A900}  % markers, alerts
\definecolor{light-bg}{HTML}{E8EDF5}          % title / section backgrounds
\definecolor{jet}{HTML}{1A1A1A}               % body text

% Semantic colors (used in diagrams and annotations)
\definecolor{positive}{HTML}{15803D}          % good / observed / identified
\definecolor{negative}{HTML}{B91C1C}          % bad / problematic / bias
\definecolor{neutral}{HTML}{525252}           % reference / context

% Highlight accents (match .hi-* classes in the SCSS)
\definecolor{hi-slate}{HTML}{314F4F}
\definecolor{hi-green}{HTML}{2E8B57}
\definecolor{hi-red}{HTML}{C41E3A}

% Semantic aliases so skill templates and snippets can write
% \textcolor{accent}{...} without hard-coding "primary-blue".
\colorlet{accent}{primary-blue}
\colorlet{accent2}{primary-gold}

% -----------------------------------------------------------------------------
% Hyperref — loaded AFTER xcolor + palette so link colors resolve.
% -----------------------------------------------------------------------------
\usepackage{hyperref}
\hypersetup{colorlinks=true, linkcolor=primary-blue, urlcolor=primary-blue,
            citecolor=primary-blue, pdfborder={0 0 0}}

% -----------------------------------------------------------------------------
% Course code — set once per deck (after \input{header}, before \begin{document})
% via \coursecode{CS401}. Shown in the footer on every slide so a deck's course
% is identifiable without opening the title page. Empty by default.
% -----------------------------------------------------------------------------
\makeatletter
\newcommand{\coursecode}[1]{\gdef\@coursecodetext{#1}}
\gdef\@coursecodetext{}
\makeatother

% -----------------------------------------------------------------------------
% Beamer theme assignments (only apply under Beamer).
%
% We detect Beamer via \@ifclassloaded{beamer}, which is the documented,
% non-internal way. This must be wrapped in \makeatletter/\makeatother
% because \@ifclassloaded contains an @-catcode character.
% -----------------------------------------------------------------------------
\makeatletter
\@ifclassloaded{beamer}{%
  \usefonttheme{professionalfonts}
  \setbeamercolor{structure}{fg=primary-blue}
  \setbeamercolor{frametitle}{fg=primary-blue}
  \setbeamercolor{title}{fg=primary-blue}
  \setbeamercolor{subtitle}{fg=primary-gold}
  \setbeamercolor{author}{fg=primary-blue}
  \setbeamercolor{institute}{fg=neutral}
  \setbeamercolor{date}{fg=primary-gold}
  \setbeamercolor{itemize item}{fg=primary-blue}
  \setbeamercolor{itemize subitem}{fg=primary-gold}
  \setbeamercolor{alerted text}{fg=primary-gold}
  \setbeamercolor{block title}{fg=white, bg=primary-blue}
  \setbeamercolor{block body}{bg=light-bg}
  % Semantic block variants: block=definition/concept (blue), exampleblock=
  % worked example (green/positive), alertblock=key takeaway or caution (gold).
  % Keep to <=2 colored boxes per slide across ALL three variants (INV-7).
  \setbeamercolor{block title example}{fg=white, bg=positive}
  \setbeamercolor{block body example}{bg=light-bg}
  \setbeamercolor{block title alerted}{fg=white, bg=primary-gold}
  \setbeamercolor{block body alerted}{bg=light-bg}

  % Minimal footer: course code (left) + page X / N (right), no navigation clutter
  \setbeamertemplate{navigation symbols}{}
  \setbeamertemplate{footline}{%
    \color{neutral}\footnotesize\hspace{0.7em}\@coursecodetext\hfill
    \insertframenumber/\inserttotalframenumber\hspace{0.7em}\vspace{0.3em}}
}{}
\makeatother

% -----------------------------------------------------------------------------
% TikZ setup — libraries the snippet gallery uses
% -----------------------------------------------------------------------------
\usepackage{tikz}
\usetikzlibrary{arrows.meta, positioning, calc, decorations.pathreplacing,
                fit, shapes.geometric, backgrounds}

% Shared TikZ styles — snippets and hand-written diagrams can pick these up
% via `\begin{tikzpicture}[every node/.style={...}]` or by naming specific
% styles (e.g., `\node[dag-node] (X) at (0,0) {X};`).
\tikzset{
  % Node styles for causal DAGs and similar
  dag-node/.style = {
    draw=primary-blue, fill=light-bg, rounded corners=2pt,
    minimum width=1.2cm, minimum height=0.9cm, align=center, font=\small
  },
  decision-node/.style = {
    draw=primary-gold, fill=white, diamond, aspect=1.8,
    minimum width=1.8cm, minimum height=1.2cm, align=center, font=\small,
    inner sep=1pt
  },
  % Edge styles — observed vs counterfactual
  observed-edge/.style = {-{Latex[length=2.5mm]}, thick, draw=jet},
  counterfactual-edge/.style = {-{Latex[length=2.5mm]}, thick, draw=neutral, dashed},
  confound-edge/.style = {-{Latex[length=2.5mm]}, thick, draw=negative, dashed},
  % Data-point styles for plots
  observed-dot/.style = {fill=primary-blue, draw=primary-blue, circle, inner sep=1.2pt},
  counterfactual-dot/.style = {fill=white, draw=neutral, circle, inner sep=1.2pt}
}

% -----------------------------------------------------------------------------
% Convenience macros
% -----------------------------------------------------------------------------
\newcommand{\muted}[1]{\textcolor{neutral}{#1}}
\newcommand{\key}[1]{\textcolor{primary-gold}{\textbf{#1}}}
\newcommand{\good}[1]{\textcolor{positive}{#1}}
\newcommand{\bad}[1]{\textcolor{negative}{#1}}

% Standout frame macro: full-bleed dark background for section transitions.
% Beamer-only (uses \begin{frame}/\setbeamercolor/\usebeamerfont) -- do not
% call from an article-class file (e.g. Notes/<CODE>/*-notes.tex). \muted,
% \key, \good, \bad, and \coursecode above are plain xcolor/gdef and are
% safe to use from either class; verified when Notes/article-class support
% was added.
% Expand: \transitionslide{Your transition title}
\newcommand{\transitionslide}[1]{%
  \begingroup
  \setbeamercolor{background canvas}{bg=primary-blue}
  \begin{frame}[plain]
    \centering
    \vfill
    {\usebeamerfont{title}\color{white}\Large #1\par}
    \vfill
  \end{frame}
  \endgroup
}

% Section divider frame (Beamer-only), an alternative to \transitionslide with a
% darker two-line style: near-black (jet) background, a small white label line
% above a larger gold title, and a thin gold rule along the bottom edge.
% Expand: \sectiondivider{Section 2}{Pointers}
\newcommand{\sectiondivider}[2]{%
  \begingroup
  \setbeamercolor{background canvas}{bg=jet}
  \begin{frame}[plain]
    \vspace*{3.0cm}
    {\color{white}\ttfamily\large #1\par}
    \vspace{0.5cm}
    {\color{primary-gold}\LARGE #2\par}
    \begin{tikzpicture}[remember picture, overlay]
      \draw[primary-gold, line width=2.5pt]
        ($(current page.south west)+(0,0.1)$) -- ($(current page.south east)+(0,0.1)$);
    \end{tikzpicture}
  \end{frame}
  \endgroup
}

% -----------------------------------------------------------------------------
% Channel watermark — "Concepts That Click" (YouTube).
%
% Article-class files (CompetitiveExam Q/A sets, Notes, Assignments) get a
% light diagonal draft watermark on every page plus a centered footer naming
% the channel. Beamer decks get a subtle TikZ background watermark instead
% (the footline already carries the course code). Centralized here so every
% MCQ PDF is branded without per-file edits.
% -----------------------------------------------------------------------------
\makeatletter
\@ifclassloaded{article}{%
  \usepackage{draftwatermark}
  \SetWatermarkText{Concepts That Click}
  \SetWatermarkScale{1}
  \SetWatermarkFontSize{2cm}
  \SetWatermarkLightness{0.86}
  \SetWatermarkAngle{45}
  \usepackage{fancyhdr}
  \pagestyle{fancy}
  \fancyhf{}
  \fancyfoot[C]{\footnotesize\textcolor{neutral}{Concepts That Click~|~YouTube: \href{https://www.youtube.com/@concepts.thatclick}{@concepts.thatclick}}}
  \renewcommand{\headrulewidth}{0pt}
  \renewcommand{\footrulewidth}{0pt}
  \fancypagestyle{plain}{%
    \fancyhf{}%
    \fancyfoot[C]{\footnotesize\textcolor{neutral}{Concepts That Click~|~YouTube: \href{https://www.youtube.com/@concepts.thatclick}{@concepts.thatclick}}}%
    \renewcommand{\headrulewidth}{0pt}%
    \renewcommand{\footrulewidth}{0pt}%
  }
}{}
\@ifclassloaded{beamer}{%
  \setbeamertemplate{background}{%
    \begin{tikzpicture}[remember picture, overlay]
      \node[rotate=45, opacity=0.055, align=center]
        at (current page.center)
        {\Huge\textcolor{neutral}{Concepts That Click}};
    \end{tikzpicture}%
  }
}{}
\makeatother 
\coursecode{JKSSB}
\title{Lesson 01 Practice: What Is a Computer?}
\author{JKSSB Computer Awareness}
\date{\today}
\begin{document}
\maketitle
\noindent\textbf{Note:} Questions have been prepared with care; please
double-check the answer key and explanations for correctness.

\begin{enumerate}[label=\textbf{Q\arabic*.}, leftmargin=*]
\item \textit{[Original]} Which description best captures the basic role of
  a computer?
  \begin{enumerate}[label=\Alph*.]
    \item It accepts data, processes it under instructions, produces output,
      and can store data or results.
    \item It decides whether every input is truthful before processing.
    \item It is a display device that presents information but cannot store.
    \item It stores data but cannot process instructions.
  \end{enumerate}
\item \textit{[Original]} A clerk enters three raw marks into a spreadsheet.
  Before any total or average is calculated, the marks are best described as:
  \begin{enumerate}[label=\Alph*.]
    \item Information
    \item Data
    \item Output
    \item Storage
  \end{enumerate}
\item \textit{[Original]} Which operation is processing in a marks
  spreadsheet?
  \begin{enumerate}[label=\Alph*.]
    \item Typing the marks
    \item Displaying the average
    \item Applying a formula to calculate the average
    \item Saving the workbook
  \end{enumerate}
\item \textit{[Original]} In IPO, the letter P stands for:
  \begin{enumerate}[label=\Alph*.]
    \item Printing
    \item Processing
    \item Programming
    \item Peripheral
  \end{enumerate}
\item \textit{[Original]} Which sequence gives the IPO cycle in its usual
  order?
  \begin{enumerate}[label=\Alph*.]
    \item Output, input, processing
    \item Input, processing, output
    \item Processing, output, input
    \item Storage, input, output
  \end{enumerate}
\item \textit{[Original]} When storage is added to IPO, the four-part
  shorthand is:
  \begin{enumerate}[label=\Alph*.]
    \item IPOS
    \item PIOS
    \item OIPS
    \item SIPO
  \end{enumerate}
\item \textit{[Original]} A saved spreadsheet is an example of which
  function?
  \begin{enumerate}[label=\Alph*.]
    \item Input
    \item Processing
    \item Output
    \item Storage
  \end{enumerate}
\item \textit{[Original]} Which statement about computer accuracy is
  \textbf{NOT} correct?
  \begin{enumerate}[label=\Alph*.]
    \item Correct instructions can produce consistent calculations.
    \item A wrong formula can produce a wrong result.
    \item High processing speed proves that input data are true.
    \item Incorrect input can lead to an incorrect result.
  \end{enumerate}
\item \textit{[Original]} The phrase ``Garbage In, Garbage Out'' warns that:
  \begin{enumerate}[label=\Alph*.]
    \item Poor input can lead to poor output.
    \item Output is always more reliable than input.
    \item Storage automatically repairs incorrect data.
    \item A computer cannot process numerical data.
  \end{enumerate}
\item \textit{[Original]} Which characteristic refers to performing many
  operations quickly?
  \begin{enumerate}[label=\Alph*.]
    \item Diligence
    \item Speed
    \item Versatility
    \item Storage
  \end{enumerate}
\item \textit{[Original]} A computer repeats a long calculation without
  tiring in the human sense. This characteristic is:
  \begin{enumerate}[label=\Alph*.]
    \item Diligence
    \item Accuracy
    \item Versatility
    \item Input
  \end{enumerate}
\item \textit{[Original]} Running a word processor, spreadsheet, and browser
  on the same computer illustrates:
  \begin{enumerate}[label=\Alph*.]
    \item Diligence
    \item Versatility
    \item Storage
    \item Output
  \end{enumerate}
\item \textit{[Original]} A programmed task continues automatically after a
  user starts it and supplies the required input. This is:
  \begin{enumerate}[label=\Alph*.]
    \item Automation
    \item Data
    \item Output
    \item GIGO
  \end{enumerate}
\item \textit{[Original]} Evaluate the statements:
  \begin{enumerate}[label=\Roman*.]
    \item Data can be raw facts or symbols.
    \item Information is data organized or processed for a purpose.
    \item The same value can be data in one task and information in another.
  \end{enumerate}
  \begin{enumerate}[label=\Alph*.]
    \item I only
    \item I and II only
    \item II and III only
    \item I, II, and III
  \end{enumerate}
\item \textit{[Original]} Which is an example of computer use in
  healthcare?
  \begin{enumerate}[label=\Alph*.]
    \item Maintaining patient records and scheduling appointments
    \item Replacing every clinical decision with an automatic guarantee
    \item Turning a monitor into secondary storage
    \item Making input data truthful by processing them quickly
  \end{enumerate}
\item \textit{[Original]} Which statement is \textbf{NOT} a limitation of a
  computer as described in basic computer fundamentals?
  \begin{enumerate}[label=\Alph*.]
    \item It does not independently determine whether input is truthful.
    \item Its output may need interpretation in context.
    \item Incorrect instructions may produce incorrect results.
    \item It performs calculations quickly when given suitable instructions.
  \end{enumerate}
\item \textit{[Original]} A displayed average is output in the current task.
  If the value is typed into a new program tomorrow, it is then acting as:
  \begin{enumerate}[label=\Alph*.]
    \item Input
    \item Processing
    \item Storage
    \item Accuracy
  \end{enumerate}
\item \textit{[Original]} A user enters expense amounts and a formula totals
  them. Which pair is matched correctly?
  \begin{enumerate}[label=\Alph*.]
    \item Amounts -- output; formula -- input
    \item Amounts -- input; formula -- processing
    \item Amounts -- storage; formula -- output
    \item Amounts -- processing; formula -- storage
  \end{enumerate}
\item \textit{[Original]} Which statement about a computer's output is
  most accurate?
  \begin{enumerate}[label=\Alph*.]
    \item Output is guaranteed to be fair if the computer is fast.
    \item Output is valid even when the input is incomplete.
    \item Output depends on the input and instructions used.
    \item Output can only be shown on a monitor.
  \end{enumerate}
\item \textit{[Original]} A user types 18 instead of 81, but the spreadsheet
  correctly calculates the average from the entered values. What best
  explains the misleading result?
  \begin{enumerate}[label=\Alph*.]
    \item The formula is processing the supplied data; the input is wrong.
    \item The monitor has converted output into input.
    \item Storage has increased the marks.
    \item The computer's speed has validated the wrong entry.
  \end{enumerate}
\end{enumerate}
\end{document}
